\section{Use a closed-loop check}
\begin{center}
\begin{tikzpicture}[x=1cm,y=1cm,>=Latex,font=\scriptsize]
\node[draw,rounded corners,minimum width=2.45cm,minimum height=.95cm,
  align=center] (n0) at (0,0) {\textbf{Observe}\\water, weather, wear};
\node[draw,rounded corners,minimum width=2.45cm,minimum height=.95cm,
  align=center] (n3) at (3,0) {\textbf{Choose}\\one clear next move};
\node[draw,rounded corners,minimum width=2.45cm,minimum height=.95cm,
  align=center] (n6) at (6,0) {\textbf{Test}\\three careful attempts};
\node[draw,rounded corners,minimum width=2.45cm,minimum height=.95cm,
  align=center] (n9) at (9,0) {\textbf{Prove}\\result or no result};
\draw[->,very thick] (n0)--(n3);
\draw[->,very thick] (n3)--(n6);
\draw[->,very thick] (n6)--(n9);
\draw[->,very thick] (n9.south) -- +(0,-.5) -| (n0.south);
\end{tikzpicture}
\end{center}

\begin{tabularx}{\linewidth}{@{}p{.22\linewidth}LL@{}}
\toprule
\textbf{Question} & \textbf{Check} & \textbf{Recovery}\\
\midrule
What changed? & Name one visible or measurable fact: wind direction, depth,
water temperature, abrasion, knot shape, lure track, fish reflex, or elapsed
time. & If you cannot name the change, observe for one minute before changing
equipment.\\
What am I testing? & Change one variable only and state what success will look
like before the attempt. & Restore the last known setup; avoid changing lure,
depth, speed, and location together.\\
Is the system sound? & Trace the full chain from hand to fish: line, knot,
leader, hook, drag, handling tool, route home. & Stop on any rough line, crossed
wrap, open snap, dull point, missing tool, or unsafe condition; repair or leave.\\
What proves completion? & A repeatable result and a final tactile/visual check,
not merely finishing the written steps. & Repeat slowly under good light. If
the check is ambiguous, treat it as a failure.\\
\bottomrule
\end{tabularx}

\begin{telospanel}{Final proof before the next cast}
Say the plan in one sentence. Tug-test the terminal connection, touch-check the
last two feet of line, confirm hook point and snap closure, set the drag, and
identify the safe landing and release route. If weather, visibility, footing,
or fish condition makes that route doubtful, change the plan before casting.
\end{telospanel}
